\documentclass[11pt,a4paper]{article}

\usepackage[margin=2.5cm]{geometry}
\usepackage{booktabs}
\usepackage{microtype}
\usepackage[colorlinks=true,linkcolor=black,urlcolor=blue]{hyperref}
\usepackage{amsmath}
\usepackage{amssymb}
\usepackage{xcolor}
\usepackage{caption}
\usepackage{pifont}
\usepackage{parskip}
\usepackage{graphicx}
\usepackage{textcomp}

\title{\textbf{Handout 01}\\[0.3em]
  \large Routing Algorithm Benchmarks --- Poland \& Europe}
\author{Eduard Tykhoniuk}
\date{2026-06-13}

\begin{document}

\maketitle
\tableofcontents
\newpage

\section{Introduction}

The practical goal of the work was to build a small C++ routing-algorithm
library and benchmark its one-to-one shortest-path algorithms on real-world road
graphs imported from Geofabrik OpenStreetMap extracts.  The code was generated
by AI.  My role was to review pull requests, check that the implementations
matched the algorithms described in the handouts.

The main comparison uses the Poland graph, where all implemented algorithms fit
within the intended memory limit.  The initial intend was to compare algorithms
on the Europe graph, however, due to the lack of RAM, only some of the algorithms
were tested, more detailed in the \hyperref[sec:europe]{Section~\ref*{sec:europe}}.

All algorithms solve point-to-point shortest paths on directed weighted road
graphs.  The objective minimizes driving distance.  The implementations compute
only the shortest-path distance value; they do not reconstruct or output the
actual sequence of vertices or edges on the path.  For imported OSM edges,
the edge length is computed from the two endpoint coordinates using Haversine
distance, scaled by 1000, rounded up with \texttt{ceil}, and stored as
\texttt{uint32\_t}.  A* uses the same coordinate model as an admissible lower
bound but rounds down with \texttt{floor}.  Integer arithmetic is faster than
floating point and avoids floating-point nondeterminism.

\section{Experimental Setup}
\label{sec:setup}

\begin{itemize}
  \item CPU: AMD Ryzen 7 5700X, 8 cores / 16 threads, max boost 4.9\,GHz.
  \item RAM: 31\,GiB physical, 99\,GiB swap.
  \item Swap backing device: NVMe at 2.7\,GB/s.
  \item OS: Linux 7.0.10-arch1-1.
\end{itemize}

The intended benchmark memory policy was a 24\,GiB RAM cap for algorithm-owned
working sets.  Swap was avoided to prevent swap I/O from distorting benchmark
timings.

\begin{table}[h]
\centering
\caption{Graph statistics.}
\begin{tabular}{lrrrl}
  \toprule
  \textbf{Graph} & \textbf{Vertices} & \textbf{Directed edges} & \textbf{File size} & \textbf{Source} \\
  \midrule
  Poland  & 17,058,998  & 33,674,745  & 675\,MB  & Geofabrik, extract 2026-05-27 \\
  Europe  & 370,316,072 & 731,016,367 & 14.7\,GB & Geofabrik, extract 2026-05-27 \\
  \bottomrule
\end{tabular}
\end{table}

For the Europe extract, the node map did not fit in RAM, so a custom parser was
written that processes the data across multiple files and merges them into a
single graph.

\subsection*{Query Methodology}

Two query types were used.

\textbf{Ranged-Dijkstra} forms the main comparison.  A Dijkstra run starts from a
random source, settles up to $N_{\max}$ vertices, and samples the target from
settled ranks $[N_{\min}, N_{\max}]$; this selects medium-distance routes and
stores reference distances for validation.

\begin{itemize}
  \item Query count: 1\,000, seed~1
  \item Settled range: $N_{\min}=10^5$, $N_{\max}=10^6$
\end{itemize}

\textbf{Random pairs} use uniformly random source--target pairs with no distance
constraint and no stored reference distance.  They appear in a separate section.

\begin{itemize}
  \item Query count: 1\,000, seed~2
\end{itemize}

Each replay benchmark runs one algorithm per process.  CH-dependent algorithms
load the prebuilt Poland CH artifact where possible, so dependency load/build
time is reported separately from algorithm-specific preprocessing.

\section{Used Algorithms}
\label{sec:algorithms}

The implemented algorithms used in the report are: Dijkstra, A*, ALT,
Bidirectional Dijkstra, Bidirectional A*, Contraction Hierarchy (CH),
Arc Flags, CHASE, and Hub Labels (HL).

\begin{description}
  \item[A* and Bidirectional A*]
    Haversine lower bound:
    $h(v,t)=\lfloor \mathrm{haversine\_m}(v,t)\times1000\rfloor$.

  \item[ALT]
    Main benchmark uses 16 farthest-point landmarks with fixed seed and 4 active
    landmarks per query.  Table~\ref{tab:alt-strategy} compares farthest, random,
    and planar selection strategies.

  \item[Arc Flags]
    Main row uses 32 Inertial regions and 16 preprocessing threads.  Flags are
    stored in a 64-bit region mask, so unsupported region counts above 64 are
    rejected rather than truncated.

  \item[CHASE]
    Uses 64 Inertial regions and a 5\% CH core fraction.

  \item[Hub Labels]
    \texttt{label\_fraction=0.05} due to high memory usage ($\approx$ 24 GiB peak).
\end{description}

\section{Poland Results}
\label{sec:poland}

All main Poland rows use 1\,000 ranged queries with seed~1.  The generated tables
below are produced directly from result JSON files, and the exact filenames are
listed at the end of the report.

% Auto-generated by LaTex/generate_tables.py -- do not edit by hand

\subsection*{Preprocessing}

\begin{table}[ht]
\centering
\caption{Preprocessing times and peak RSS --- Poland graph.}
\label{tab:preprocess-poland}
\resizebox{\textwidth}{!}{\begin{tabular}{lrrrrrl}
  \toprule
  \textbf{Algorithm} & \textbf{Dep.\ wall\,(s)} & \textbf{Dep.\ CPU\,(s)} & \textbf{Algo\ wall\,(s)} & \textbf{Algo\ CPU\,(s)} & \textbf{RSS\ after} & \textbf{Notes} \\
  \midrule
  Dijkstra & $<\!0.001$ & $<\!0.001$ & $<\!0.001$ & $<\!0.001$ & 844\,MB & --- \\
  A* & $<\!0.001$ & $<\!0.001$ & $<\!0.001$ & $<\!0.001$ & 1.1\,GB & --- \\
  Bidirectional Dijkstra & $<\!0.001$ & $<\!0.001$ & 1.0 & 0.9 & 1.8\,GB & --- \\
  Bidirectional A* & $<\!0.001$ & $<\!0.001$ & 1.0 & 0.9 & 2.0\,GB & --- \\
  ALT (16 landmarks) & $<\!0.001$ & $<\!0.001$ & 13.3 & 13.2 & 6.2\,GB & --- \\
  CH & $<\!0.001$ & $<\!0.001$ & 425.0 & 422.8 & 8.7\,GB & CH saved, +64M aux.~edges \\
  Arc Flags (32 regions) & 2.5 & 2.4 & 1\,984.6 & 29\,978.0 & 6.9\,GB & CH loaded \\
  CHASE (64 regions) & 2.3 & 2.3 & 16.2 & 16.1 & 3.4\,GB & CH loaded \\
  Hub Labels (0.05) & 2.3 & 2.2 & 397.4 & 394.5 & 24.0\,GB & CH loaded \\
  \bottomrule
\end{tabular}}
\end{table}

\subsection*{Query performance}

\begin{table}[ht]
\centering
\caption{Query latency --- Poland graph, 1\,000 queries, seed~1, settled $\in[10^5,10^6]$.}
\label{tab:query-poland}
\resizebox{\textwidth}{!}{\begin{tabular}{lrrrrrrl}
  \toprule
  \textbf{Algorithm} & \textbf{Mean\,({\textmu}s)} & \textbf{P50} & \textbf{P95} & \textbf{P99} & \textbf{Mean\ CPU\,({\textmu}s)} & \textbf{Speedup$^\dagger$} & \textbf{Validated} \\
  \midrule
  Dijkstra & 54\,379 & 53\,304 & 99\,190 & 107\,870 & 54\,102 & 1.0$\times$ & $\checkmark$ \\
  A* & 21\,027 & 16\,370 & 54\,021 & 74\,173 & 20\,921 & 2.6$\times$ & $\checkmark$ \\
  Bidirectional Dijkstra & 44\,100 & 38\,290 & 104\,369 & 136\,532 & 43\,871 & 1.2$\times$ & $\checkmark$ \\
  Bidirectional A* & 49\,435 & 38\,825 & 125\,996 & 193\,506 & 49\,184 & 1.1$\times$ & $\checkmark$ \\
  ALT (16 landmarks) & 24\,252 & 17\,361 & 65\,848 & 83\,936 & 24\,127 & 2.2$\times$ & $\checkmark$ \\
  CH & 169.0 & 158.7 & 318.9 & 392.7 & 168.4 & 322$\times$ & $\checkmark$ \\
  Arc Flags (32 regions) & 20\,011 & 17\,405 & 43\,948 & 50\,709 & 19\,910 & 2.7$\times$ & $\checkmark$ \\
  CHASE (64 regions) & 179.9 & 164.7 & 341.3 & 423.0 & 179.2 & 302$\times$ & $\checkmark$ \\
  Hub Labels (0.05) & 78.5 & 78.4 & 111.7 & 132.2 & 78.1 & 693$\times$ & $\checkmark$ \\
  \bottomrule
\end{tabular}}
\noindent\footnotesize $^\dagger$Speedup = Dijkstra mean wall / algorithm mean wall.\\
\end{table}



% Auto-generated by LaTex/generate_tables.py -- do not edit by hand
\begin{table}[ht]
\centering
\caption{Arc Flags configuration sweep --- Poland graph, 1\,000 ranged queries, seed~1.}
\label{tab:arcflags-sweep}
\resizebox{\textwidth}{!}{%
\begin{tabular}{lrrrrrrrr}
  \toprule
  \textbf{Variant} & \textbf{Threads} & \textbf{Pre.\ wall} & \textbf{Pre.\ CPU} & \textbf{RSS} & \textbf{Mean\,(\textmu{}s)} & \textbf{P99\,(\textmu{}s)} & \textbf{Prune rate} & \textbf{Speedup} \\
  \midrule
  Inertial-32 & 1 & 71.4\,min & 4\,253\,s & 4.9\,GB & 21\,668 & 57\,648 & 2.7\% & 2.5$\times$ \\
  Inertial-32 & 16 & 33.1\,min & 29\,978\,s & 6.9\,GB & 20\,011 & 50\,709 & 2.7\% & 2.7$\times$ \\
  Inertial-64 & 16 & 49.9\,min & 41\,951\,s & 6.9\,GB & 14\,997 & 50\,176 & 4.3\% & 3.6$\times$ \\
  KaMinPar-64 & 16 & 8.2\,min & 7\,215\,s & 7.2\,GB & 12\,410 & 29\,331 & 4.5\% & 4.4$\times$ \\
  \bottomrule
\end{tabular}}
\smallskip
\noindent\footnotesize Prune rate = mean pruned arcs / mean relaxed arcs per query.\\

\end{table}



% Auto-generated by LaTex/generate_tables.py -- do not edit by hand
\begin{table}[ht]
\centering
\caption{ALT landmark-selection strategies --- Poland graph, 1\,000 ranged queries, seed~1.}
\label{tab:alt-strategy}
\resizebox{\textwidth}{!}{%
\begin{tabular}{lrrrrr}
  \toprule
  \textbf{Variant} & \textbf{Pre.\ wall\,(s)} & \textbf{RSS} & \textbf{Mean\,(\textmu{}s)} & \textbf{P99\,(\textmu{}s)} & \textbf{Speedup} \\
  \midrule
  farthest-16-active-4 & 13.3 & 6.2\,GB & 24\,252 & 83\,936 & 2.2$\times$ \\
  random-16-active-4 & 53.3 & 6.2\,GB & 8\,601 & 47\,667 & 6.3$\times$ \\
  planar-16-active-4 & 24.6 & 6.2\,GB & 6\,293 & 31\,116 & 8.6$\times$ \\
  \bottomrule
\end{tabular}}
\end{table}



Hub Labels is the fastest measured algorithm on Poland.  After HL, CH is the
fastest full-coverage CH-family baseline in this implementation.  CHASE is close
to CH but does not improve query time because CH already leaves a tiny search
space for the extra flag checks to prune.

The unfortunate result for Arc Flags is that substantial preprocessing produces
only a slight query speedup.  Inertial partitioning is geometric rather than
graph-aware, so many edges cross region boundaries and carry multiple region
bits.  The 1-to-16 thread preprocessing change mostly exposes memory-bandwidth
contention over large PHAST and flag arrays.

\section{Random-Pair Queries}
\label{sec:random-pairs}

All Poland algorithms were also benchmarked on 1\,000 uniformly random
source--target pairs (seed~2).  Random pairs do not carry stored reference
distances, so results are timing-only.

% Auto-generated by LaTex/generate_tables.py -- do not edit by hand
\begin{table}[ht]
\centering
\caption{Random-pair check --- Poland graph, 1\,000 random pairs, seed~2.}
\label{tab:random-pairs}
\resizebox{\textwidth}{!}{\begin{tabular}{lrrrrrrl}
  \toprule
  \textbf{Algorithm} & \textbf{Mean\,({\textmu}s)} & \textbf{P50} & \textbf{P95} & \textbf{P99} & \textbf{Mean\ CPU\,({\textmu}s)} & \textbf{Speedup$^\dagger$} & \textbf{Validated} \\
  \midrule
  Dijkstra & 854\,241 & 848\,338 & 1\,667\,118 & 1\,751\,287 & 850\,134 & 1.0$\times$ & --- \\
  A* & 458\,316 & 320\,566 & 1\,223\,368 & 4\,023\,714 & 456\,140 & 1.9$\times$ & --- \\
  Bidirectional Dijkstra & 680\,702 & 637\,070 & 1\,488\,858 & 1\,682\,276 & 677\,420 & 1.3$\times$ & --- \\
  Bidirectional A* & 918\,029 & 699\,416 & 2\,573\,625 & 3\,813\,021 & 913\,639 & 0.9$\times$ & --- \\
  ALT (farthest, 16/4) & 456\,801 & 310\,669 & 1\,266\,264 & 2\,716\,578 & 454\,591 & 1.9$\times$ & --- \\
  CH & 695.5 & 744.8 & 1\,104 & 1\,209 & 692.2 & 1\,228$\times$ & --- \\
  Arc Flags (KaMinPar-64) & 31\,489 & 28\,266 & 66\,801 & 108\,520 & 31\,327 & 27.1$\times$ & --- \\
  CHASE (64 regions) & 797.0 & 852.6 & 1\,317 & 1\,485 & 792.6 & 1\,072$\times$ & --- \\
  Hub Labels (0.05) & 71.1 & 72.3 & 100.4 & 112.1 & 70.8 & 12\,012$\times$ & --- \\
  \bottomrule
\end{tabular}}
\noindent\footnotesize Random pairs do not carry reference distances; \texttt{distance\_validated=false}.\\
\end{table}




\section{Europe Results}
\label{sec:europe}

The Europe graph is large enough that memory behavior dominates the feasible
benchmark set.  Dijkstra, A*, and Bidirectional Dijkstra were run; ALT, CH,
Arc Flags, CHASE, Hub Labels, and Bidirectional A* exceed the intended RAM cap.

% Auto-generated by LaTex/generate_tables.py -- do not edit by hand

\subsection*{Preprocessing}

\begin{table}[ht]
\centering
\caption{Preprocessing times and peak RSS --- Europe graph.}
\label{tab:preprocess-europe}
\resizebox{\textwidth}{!}{\begin{tabular}{lrrrrrl}
  \toprule
  \textbf{Algorithm} & \textbf{Dep.\ wall\,(s)} & \textbf{Dep.\ CPU\,(s)} & \textbf{Algo\ wall\,(s)} & \textbf{Algo\ CPU\,(s)} & \textbf{RSS\ after} & \textbf{Notes} \\
  \midrule
  Dijkstra & $<\!0.001$ & $<\!0.001$ & $<\!0.001$ & $<\!0.001$ & 17.8\,GB & --- \\
  A* & $<\!0.001$ & $<\!0.001$ & $<\!0.001$ & $<\!0.001$ & 23.3\,GB & --- \\
  Bidirectional Dijkstra & $<\!0.001$ & $<\!0.001$ & 116.9 & 95.9 & 28.4\,GB & --- \\
  \bottomrule
\end{tabular}}
\end{table}

\subsection*{Query performance}

\begin{table}[ht]
\centering
\caption{Query latency --- Europe graph, 1\,000 queries, seed~1, settled $\in[10^5,10^6]$.}
\label{tab:query-europe}
\resizebox{\textwidth}{!}{\begin{tabular}{lrrrrrrl}
  \toprule
  \textbf{Algorithm} & \textbf{Mean\,({\textmu}s)} & \textbf{P50} & \textbf{P95} & \textbf{P99} & \textbf{Mean\ CPU\,({\textmu}s)} & \textbf{Speedup$^\dagger$} & \textbf{Validated} \\
  \midrule
  Dijkstra & 63\,905 & 61\,885 & 122\,901 & 146\,320 & 63\,589 & 1.0$\times$ & $\checkmark$ \\
  A* & 25\,129 & 20\,098 & 63\,040 & 93\,098 & 25\,001 & 2.5$\times$ & $\checkmark$ \\
  Bidirectional Dijkstra & 212\,856 & 172\,723 & 530\,239 & 919\,133 & 121\,590 & 0.3$\times$ & $\checkmark$ \\
  \bottomrule
\end{tabular}}
\noindent\footnotesize $^\dagger$Speedup = Dijkstra mean wall / algorithm mean wall.\\
\end{table}



A* works well at this big scale because the Haversine heuristic is a very good approximation
for long geographic routes.  It needs the graph, coordinate sidecar, and scratch
arrays, but it does not need and preprocessing.  In contrast, ALT would
need roughly 47\,GB for 16 forward/reverse landmark tables on 370\,M vertices,
and CH-family preprocessing is far beyond the available RAM.

\section{Discussion}
\label{sec:discussion}

The main comparison is useful but not a final best-possible performance study:
the algorithms were implemented functionally and with attention to obvious hot
paths, but not exhaustively optimized.  That caveat matters especially for
Arc Flags and Hub Labels, where implementation details can strongly affect
preprocessing memory and query latency.

On Poland, CH and Hub Labels dominate query latency because their preprocessing
compresses the search space.  A* is the best low-preprocessing option: it
relies only on coordinates and scratch memory.  ALT improves over Dijkstra but
pays a dense landmark-table cost.  Arc Flags benefits from a graph-aware
partition: \href{https://github.com/KaHIP/KaMinPar}{KaMinPar}-64
(Table~\ref{tab:arcflags-sweep}) improves over inertial partitioning.

CH class family of algorithms worked very well, as they greatly compressed the number of vertices
in the optimal path.
This also may be confirmed by the fact, that speedup is higher on completely random pairs as the search space
does not rely on the underlying (real) shortest path's length.

However, this family requires long preprocessing time as well as a high memory usage.
On this hardware, none of the preprocessing algorithms worked for Europe.
The other family of algorithms is low-preprocessing algorithms: A* was the best
Europe-scale option.

\section{Failures and Learning Process}
\label{sec:failures}

CH and Arc Flags artifact saving was added because their preprocessing is the
same between benchmark runs.  Reusing those artifacts avoids rerunning the same
expensive preprocessing step and makes it much faster to get benchmark results.
Automated integration tests were also added to cross-check algorithm distances
against Dijkstra reference results.

Memory pressure was the main practical limitation.  Large benchmark batch sizes
in PHAST consumed too much RAM, so the implementation was reduced to
\texttt{kBatchSize = 1}.  Hub Labels with
\texttt{label\_fraction=0.25} also turned out to be too large for the available
memory, so the reported Poland run uses a smaller fraction.

For graph partitioning, \href{https://github.com/KaHIP/KaMinPar}{KaMinPar} was
used as a newer graph-partitioning library than METIS.  It improved over the
simple inertial partitioning variant, but did not change the broader conclusion
that Arc Flags remained preprocessing-heavy for the measured speedup.

\clearpage
% Auto-generated by LaTex/generate_tables.py -- do not edit by hand
\section*{Result JSON Files Used}
\addcontentsline{toc}{section}{Result JSON Files Used}
\small
\resizebox{\textwidth}{!}{%
\begin{tabular}{llllllr}
  \toprule
  \textbf{Table group} & \textbf{JSON file} & \textbf{Commit} & \textbf{Variant} & \textbf{Graph} & \textbf{Query set} & \textbf{Queries} \\
  \midrule
  Main Poland & dijkstra\_default\_poland\_20260613T040011.json &  & default & poland-driving-20260613.graph & poland\_ranged\_q1000\_s1.json & 1000 \\
  Main Poland & astar\_default\_poland\_20260613T040113.json &  & default & poland-driving-20260613.graph & poland\_ranged\_q1000\_s1.json & 1000 \\
  Main Poland & bidijkstra\_default\_poland\_20260613T040232.json &  & default & poland-driving-20260613.graph & poland\_ranged\_q1000\_s1.json & 1000 \\
  Main Poland & bidi\_astar\_default\_poland\_20260613T040317.json &  & default & poland-driving-20260613.graph & poland\_ranged\_q1000\_s1.json & 1000 \\
  Main Poland & alt\_farthest-16-active-4\_poland\_20260613T203137.json &  & farthest-16-active-4 & poland-driving-20260613.graph & poland\_ranged\_q1000\_s1.json & 1000 \\
  Main Poland & ch\_default\_poland\_20260613T040417.json &  & default & poland-driving-20260613.graph & poland\_ranged\_q1000\_s1.json & 1000 \\
  Main Poland & arcflags\_inertial-32\_bs1\_16t\_poland\_20260613T173814.json &  & 32-partitions-inertial & poland-driving-20260613.graph & poland\_ranged\_q1000\_s1.json & 1000 \\
  Main Poland & chase\_default\_poland\_20260613T042145.json &  & 5pct-threshold-64-partitions-inertial & poland-driving-20260613.graph & poland\_ranged\_q1000\_s1.json & 1000 \\
  Main Poland & hl\_fraction-0.05\_poland\_20260613T113450.json &  & label\_fraction=0.05, memory\_budget=22GB & poland-driving-20260613.graph & poland\_ranged\_q1000\_s1.json & 1000 \\
  Main Europe & dijkstra\_default\_europe\_20260613T051741.json &  & default & europe-driving-20260613.graph & europe\_ranged\_q1000\_s1.json & 1000 \\
  Main Europe & astar\_default\_europe\_20260613T190438.json &  & default & europe-driving-20260613.graph & europe\_ranged\_q1000\_s1.json & 1000 \\
  Main Europe & bidijkstra\_default\_europe\_20260613T052957.json &  & default & europe-driving-20260613.graph & europe\_ranged\_q1000\_s1.json & 1000 \\
  Arc Flags sweep & arcflags\_default\_poland\_20260613T041134.json &  & 32-partitions-inertial & poland-driving-20260613.graph & poland\_ranged\_q1000\_s1.json & 1000 \\
  Arc Flags sweep & arcflags\_inertial-32\_bs1\_16t\_poland\_20260613T173814.json &  & 32-partitions-inertial & poland-driving-20260613.graph & poland\_ranged\_q1000\_s1.json & 1000 \\
  Arc Flags sweep & arcflags\_inertial-64\_bs1\_16t\_poland\_20260613T181258.json &  & 64-partitions-inertial & poland-driving-20260613.graph & poland\_ranged\_q1000\_s1.json & 1000 \\
  Arc Flags sweep & arcflags\_kaminpar-64\_16t\_poland\_20260613T204500.json &  & 64-partitions-kaminpar & poland-driving-20260613.graph & poland\_ranged\_q1000\_s1.json & 1000 \\
  ALT strategy & alt\_farthest-16-active-4\_poland\_20260613T203137.json &  & farthest-16-active-4 & poland-driving-20260613.graph & poland\_ranged\_q1000\_s1.json & 1000 \\
  ALT strategy & alt\_random-16-active-4\_poland\_20260613T203227.json &  & random-16-active-4 & poland-driving-20260613.graph & poland\_ranged\_q1000\_s1.json & 1000 \\
  ALT strategy & alt\_planar-16-active-4\_poland\_20260613T203340.json &  & planar-16-active-4 & poland-driving-20260613.graph & poland\_ranged\_q1000\_s1.json & 1000 \\
  Random pairs & dijkstra\_random-pairs\_poland\_20260613T205100.json &  & random-pairs & poland-driving-20260613.graph & poland\_random\_q1000\_s2.json & 1000 \\
  Random pairs & astar\_random-pairs\_poland\_20260613T205100.json &  & random-pairs & poland-driving-20260613.graph & poland\_random\_q1000\_s2.json & 1000 \\
  Random pairs & bidijkstra\_random-pairs\_poland\_20260613T205100.json &  & random-pairs & poland-driving-20260613.graph & poland\_random\_q1000\_s2.json & 1000 \\
  Random pairs & bidi\_astar\_random-pairs\_poland\_20260613T205100.json &  & random-pairs & poland-driving-20260613.graph & poland\_random\_q1000\_s2.json & 1000 \\
  Random pairs & alt\_random-pairs-farthest-16-active-4\_poland\_20260613T205100.json &  & random-pairs-farthest-16-active-4 & poland-driving-20260613.graph & poland\_random\_q1000\_s2.json & 1000 \\
  Random pairs & ch\_random-pairs\_poland\_20260613T205100.json &  & random-pairs & poland-driving-20260613.graph & poland\_random\_q1000\_s2.json & 1000 \\
  Random pairs & arcflags\_random-pairs-kaminpar-64\_16t\_poland\_20260613T205100.json &  & random-pairs-kaminpar-64-16t & poland-driving-20260613.graph & poland\_random\_q1000\_s2.json & 1000 \\
  Random pairs & chase\_random-pairs\_poland\_20260613T205100.json &  & random-pairs & poland-driving-20260613.graph & poland\_random\_q1000\_s2.json & 1000 \\
  Random pairs & hl\_random-pairs\_fraction-0.05\_poland\_20260613T205100.json &  & random-pairs-label\_fraction=0.05 & poland-driving-20260613.graph & poland\_random\_q1000\_s2.json & 1000 \\
  \bottomrule
\end{tabular}}
\normalsize



\end{document}
